\section{Síntesis y ampliación}

Esta clase responde con cuatro temas qué significa «Beyond LLMs». Emergent abilities nos recuerda distinguir entre un cambio real de capacidad y un metric artifact; Intermediate-Guided Reasoning extiende una sola generación a chain, tree, decomposition o program execution; BabyLM usa un presupuesto de datos restringido para indagar en la sample efficiency; y Agents coloca al modelo dentro de un sistema capaz de actuar, recordar, colaborar y corregir errores.

\subsection{Marco unificado: la capacidad proviene del modelo, del proceso y del ciclo cerrado}

Las dos mitades de la clase pueden unificarse en tres capas: el model aporta representación y generación probabilísticas, el inference procedure organiza los intermediate states y el system loop conecta los resultados con el environment. Una escala mayor no produce automáticamente un razonamiento fiel, una trayectoria de razonamiento más larga no produce automáticamente una acción correcta y más herramientas no producen automáticamente un sistema confiable. Cada capa requiere métricas propias y su propio análisis de fallas.

\begin{importantbox}{Ocho conclusiones aplicables de esta sesión}
\begin{enumerate}
  \item Al reportar emergence, revisar al mismo tiempo la metric, el prompt, la task distribution y una métrica continua alternativa;
  \item Para cada reasoning method, distinguir entre final-answer accuracy, step validity y compute cost;
  \item Delegar con prioridad el cálculo exacto a un interpreter y verificar que el programa sea coherente con la semántica del problema;
  \item Al comparar modelos pequeños, reportar al mismo tiempo data budget, architecture, objective y curriculum;
  \item Evaluar por separado el foundation model y el agent system;
  \item Definir para cada acción real sus permissions, su success predicate y su confirmation gate;
  \item En Multi-agent, usar un task state estructurado, reintentos limitados y evidence aggregation;
  \item La ejecución de larga duración debe apoyarse en external feedback, audit, sandbox y recovery.
\end{enumerate}
\end{importantbox}

\subsection{Cómo verificar una Agent Claim}

Ante afirmaciones como «el agent puede reservar boletos, aprobar exámenes u operar cualquier software», puede revisarse en cinco pasos: primero, fijar la tarea y la versión; después, hacer públicos el estado inicial y los permisos; luego, ejecutar varias veces y guardar las trajectories completas; determinar el éxito con la postcondition del entorno y no con la autoevaluación del modelo; por último, reportar las categorías de falla, la intervención humana, el costo y los eventos irreversibles. Solo así una demo se convierte en evidence comparable.

\begin{warningbox}{No presentar la visión de la clase de 2023 como hechos de productos de 2026}
Esta sesión registra el contenido del curso del 2023-11-28. MultiOn, AutoGPT, Action API y LLM OS surgieron en un contexto histórico específico; las notas conservan su valor didáctico, pero no se apoyan en ellos para afirmar las capacidades actuales de los productos, la situación de las empresas ni los estándares de la industria.
\end{warningbox}

\subsection{Lecturas adicionales y resumen final}

\enlargethispage{5\baselineskip}

Se recomienda continuar por tres rutas: primero, leer el debate sobre emergent abilities y metric-choice, y comparar el exact-match discreto con la pérdida continua; segundo, leer Chain-of-Thought, Tree of Thoughts, PAL, Program of Thoughts y BabyLM Challenge para entender el proceso de razonamiento y la eficiencia de muestras; tercero, leer la revisión de Lilian Weng sobre LLM-powered autonomous agents y reexaminar la arquitectura vista en clase a la luz de la agent evaluation moderna, la computer-use safety y el diseño de sandbox. Al leer, conviene dar prioridad a los artículos, las páginas oficiales de los proyectos y los experimentos reproducibles, sin sustituir la evidencia por demos de productos.

El verdadero «Beyond LLMs» no consiste en abandonar los LLM, sino en dejar de considerar una sola generación de texto como el punto final de la capacidad. Un sistema más confiable elige la intermediate representation adecuada, obtiene resultados verificables mediante herramientas, guarda en memory estados con provenance, ejecuta bounded tasks en paralelo con multi-agent y realiza acciones en el mundo real bajo permissions, feedback, reflection, sandboxing y human oversight.

\end{document}
